%% kapa: exam=1 round=28 year=2017 period=1 form=A subject=civil_law subject_ko=민법 seq=15 number=15
%% source: Q-Net 감정평가사 기출문제 (cst00309) artlSeq=5199603 file=extracted/1008/5199603/2017년 제28회 감정평가사 1차시험 문제지/1교시 A형.pdf
%% conversion: kordoc 4.24.51 -> tools/gen1.py (auto); review=proofread-sonnet; flags=none
\begin{question}{15}
법률행위의 취소에 관한 설명으로 옳지 않은 것은? (다툼이 있으면 판례에 따름)
\begin{choices}
  \item 사기를 이유로 취소된 법률행위는 처음부터 무효인 것으로 본다.
  \item 제한능력자가 취소권을 가지는 경우 법정대리인의 동의 없이 행사할 수 있다.
  \item 피성년후견인은 법정대리인의 동의가 있더라도 재산상 법률행위를 스스로 유효하게 추인할 수 없다.
  \item 법정대리인이 미성년자의 법률행위를 추인하는 경우, 취소 원인이 소멸된 후에 하여야만 효력이 있다.
  \item 법률행위를 취소한 후라도 무효행위의 추인의 요건과 효력으로서 추인할 수 있다.
\end{choices}
\end{question}
\answer{4}
